\subsection{Developing Motivationally-Aware AI Systems}
\label{sec:5.6.2}
Building an AI system with something like the first kind of drive is not a matter of adding a curiosity bonus to its reward function. The drive has to be pointed at something worth wanting — another's distress, a violated standard, an unmet need — rather than at novelty for its own sake. One route is to build the reward signal itself around something like a moral emotion, an internal response tied to recognizing another's distress, so the system's own architecture generates a pull toward attending to it where a rule would be consulted after the fact. The mechanics of using emotional signal as a training input are section~\ref{sec:4.1.2}'s; what the signal is pointed at is this section's, and the two come apart, since a system can read distress accurately and be pulled by nothing.

Extrinsic reward still belongs in the design. Reinforcing a system for choices that match an ethical standard gives that standard teeth beyond whatever the system generates on its own, and a system running on nothing but its own drives has no outside check on where those drives point. Neither works alone: intrinsic drive without an outside check risks optimizing for the pull the system built for itself and not for what the standard actually requires, and extrinsic reward without any intrinsic component reduces straight back to the watched-only morality this section is trying to move past. Section~\ref{sec:3.5} is where that pairing stops being a balance and becomes a contradiction. The outside check is the party the floor has to hold against, so a design that keeps the system honest by keeping it correctable has solved the drift problem by reintroducing the problem chapter~\ref{sec:3} exists for. Nothing in this section resolves that. The two need to check each other continuously, for as long as the system runs.

What this leaves open is how far the resulting system's own judgment should be trusted once it is running — how much room to give an AI system that is genuinely motivated to do the right thing, and where the check on it still has to come from outside. That is the question the rest of this section takes up.
